\documentclass[11pt,letterpaper]{article}

\usepackage[spanish,es-noshorthands]{babel}
\usepackage{fontspec}
\setmainfont{Source Serif Pro}
\usepackage{microtype}
\usepackage{geometry}
\usepackage{enumitem}
\usepackage{xcolor}
\usepackage{csquotes}
\usepackage{tikz}
\usetikzlibrary{arrows.meta,positioning,fit,backgrounds}
\usepackage{xurl}
\usepackage[hidelinks]{hyperref}
\usepackage[backend=biber,style=apa,sorting=nyt]{biblatex}
\DeclareLanguageMapping{spanish}{spanish-apa}

\geometry{margin=2.3cm}
\hbadness=10000
\setlength{\emergencystretch}{1em}
\setlength{\parindent}{0pt}
\setlength{\parskip}{0.5em}
\setlist{nosep}
\addbibresource{referencias.bib}
\AtBeginBibliography{\sloppy\small}
\setlength{\bibitemsep}{0.3em}

% Colores: azul de Google Drive, verde de NotebookLM, un tono por
% departamento, gris del marco de Google Workspace y del texto secundario.
\definecolor{cDrive}{RGB}{26,95,180}
\definecolor{cNlm}{RGB}{36,122,82}
\definecolor{cVentas}{RGB}{186,86,30}
\definecolor{cOper}{RGB}{104,62,148}
\definecolor{cTec}{RGB}{26,58,102}
\definecolor{cGris}{RGB}{120,116,110}

\tikzset{
  % los rótulos de la figura no se dividen con guion
  every node/.append style={execute at begin node={\hyphenpenalty=10000}},
  depto/.style={draw=#1, fill=#1!10, very thick, rounded corners=3pt,
    text width=2.2cm, align=center, minimum height=1.05cm, font=\small\bfseries},
  carpeta/.style={draw=cDrive, fill=white, thick, rounded corners=2pt,
    text width=2.7cm, align=center, minimum height=0.85cm, font=\footnotesize},
  flecha/.style={-{Stealth[length=2.6mm]}, thick, draw=#1},
  rotulo/.style={font=\scriptsize, text=cGris!50!black, align=center},
}

\begin{document}

% --- Portada -----------------------------------------------------------------
\begin{titlepage}
\centering
\vspace*{3.2cm}
{\large\scshape Gestión del Desarrollo Tecnológico:\\
Estrategias y Análisis de Portfolio\par}
\vspace{0.5cm}
{\normalsize Módulo 4. Gestión del conocimiento y la información\par}
\vspace{2.6cm}
{\color{cTec}\rule{0.55\textwidth}{0.8pt}\par}
\vspace{0.7cm}
{\LARGE\bfseries Conocimiento compartido en una\\[0.2em]
empresa 100\% virtual\par}
\vspace{0.5cm}
{\large Cómo se crea, se guarda y se consulta el conocimiento\\
con Google Drive y NotebookLM\par}
\vspace{0.7cm}
{\color{cTec}\rule{0.55\textwidth}{0.8pt}\par}
\vfill
{\large Carlos Luis Rojas Aragonés\par}
\vspace{0.3cm}
{\normalsize Chief Technology Officer Program\par}
\vspace{0.3cm}
{\normalsize 28 de septiembre de 2026\par}
\vspace*{1.5cm}
\end{titlepage}

% --- Cuerpo ------------------------------------------------------------------
\section*{1. Cómo se gestiona y retiene el conocimiento}

En mi empresa el conocimiento nuevo entra sobre todo por las personas. Se
fomenta que los líderes técnicos asistan a conferencias de tecnología y, a su
regreso, compartan con el resto del equipo lo que se puede aplicar al negocio.
No se trata de repetir la conferencia, sino de filtrarla: quien asistió decide
qué ideas son relevantes para nosotros y las traduce a nuestro contexto.

Este paso es la conversión de conocimiento tácito en explícito que
\textcite{nonaka1995} llaman \emph{externalización}: lo que el líder aprendió
escuchando y conversando se convierte en una presentación o un documento que
otros pueden leer. Sin ese paso, el conocimiento se queda en la cabeza de
quien viajó y se pierde cuando esa persona cambia de puesto.

\section*{2. El sistema: Google Drive y NotebookLM}

Sí contamos con un sistema, y es deliberadamente sencillo. Cada departamento
(Ventas, Operaciones y Tecnología) crea y mantiene actualizada su propia
documentación en Google Drive. Sobre esos documentos usamos NotebookLM, la
herramienta de Google Workspace que responde preguntas apoyándose únicamente
en las fuentes que se le entregan y cita el fragmento de donde sale cada
respuesta \parencite{notebooklm}. La Figura~\ref{fig:ecosistema} muestra cómo
se relacionan las piezas.

\begin{figure}[h]
\centering
\begin{tikzpicture}[node distance=0.55cm]
  % Departamentos
  \node[depto=cVentas] (ven) at (0,1.55) {Ventas};
  \node[depto=cOper]   (ope) at (0,0)    {Operaciones};
  \node[depto=cTec]    (tec) at (0,-1.55) {Tecnología};

  % Unidades compartidas de Google Drive
  \node[carpeta] (dven) at (4.6,1.55)  {Unidad compartida\\ de Ventas};
  \node[carpeta] (dope) at (4.6,0)     {Unidad compartida\\ de Operaciones};
  \node[carpeta] (dtec) at (4.6,-1.55) {Unidad compartida\\ de Tecnología};
  \begin{scope}[on background layer]
    \node[draw=cDrive, fill=cDrive!8, very thick, rounded corners=4pt,
          inner sep=0.28cm, fit=(dven)(dtec), yshift=0.18cm,
          label={[font=\small\bfseries, text=cDrive, anchor=north,
                  yshift=-0.02cm]north:Google Drive}] (drive) {};
  \end{scope}

  % NotebookLM
  \node[draw=cNlm, fill=cNlm!10, very thick, rounded corners=4pt,
        text width=2.6cm, align=center, minimum height=2.3cm]
        (nlm) at (9.15,0)
        {\small\bfseries NotebookLM\\[0.3em]
         \footnotesize\mdseries un cuaderno por departamento, con las
         unidades como fuentes};

  % Personas que consultan
  \node[draw=cGris, fill=cGris!8, very thick, rounded corners=4pt,
        text width=2.2cm, align=center, minimum height=1.6cm, font=\small]
        (per) at (13.35,0)
        {\textbf{Personas}\\[0.2em]\footnotesize con años en la empresa y
         de nuevo ingreso};

  % Marco de Google Workspace
  \begin{scope}[on background layer]
    \node[draw=cGris, dashed, thick, rounded corners=6pt, inner sep=0.3cm,
          fit=(ven)(tec)(drive)(nlm)(per), yshift=0.22cm,
          label={[font=\small\bfseries, text=cGris!70!black, anchor=north west,
                  xshift=0.1cm]north west:Google Workspace}] (ws) {};
  \end{scope}

  % Conferencias, fuera del ecosistema
  \node[draw=cTec, dashed, thick, rounded corners=3pt, text width=3.2cm,
        align=center, font=\footnotesize, below=1.25cm of tec] (conf)
        {\textbf{Conferencias de tecnología}\\ a las que asisten los líderes
         técnicos};

  % Flujos
  \draw[flecha=cVentas] (ven) -- node[rotulo, above] {crea y\\[-0.2em] actualiza} (dven);
  \draw[flecha=cOper]   (ope) -- node[rotulo, above] {crea y\\[-0.2em] actualiza} (dope);
  \draw[flecha=cTec]    (tec) -- node[rotulo, above] {crea y\\[-0.2em] actualiza} (dtec);
  \draw[flecha=cTec] (conf) -- node[rotulo, right, align=left, pos=0.3]
       {comparten lo aplicable\\[-0.2em] al negocio} (tec);

  \draw[flecha=cDrive] (dven.east) -- ([yshift=0.55cm]nlm.west);
  \draw[flecha=cDrive] (dope.east) -- node[rotulo, above] {fuentes} (nlm.west);
  \draw[flecha=cDrive] (dtec.east) -- ([yshift=-0.55cm]nlm.west);

  \draw[flecha=cGris] ([yshift=0.32cm]per.west) --
       node[rotulo, above] {pregunta} ([yshift=0.32cm]nlm.east);
  \draw[flecha=cNlm] ([yshift=-0.32cm]nlm.east) --
       node[rotulo, below] {respuesta\\[-0.2em] con citas} ([yshift=-0.32cm]per.west);
\end{tikzpicture}
\caption{El ecosistema de gestión del conocimiento dentro de Google Workspace.
Cada departamento es dueño de su unidad compartida en Google Drive; NotebookLM
usa esos documentos como fuentes y responde a cualquier persona de la empresa
citando el documento de origen. Las conferencias alimentan el sistema desde
fuera, a través de los líderes técnicos.}
\label{fig:ecosistema}
\end{figure}

\section*{3. ¿Da buenos resultados?}

Creo que sí, por tres razones:

\begin{itemize}[leftmargin=1.2em, itemsep=0.25em]
  \item \textbf{Sirve a los dos extremos de la antigüedad.} A una persona nueva
    le evita tener que saber a quién preguntar; a una persona con muchos años
    en la empresa le ahorra buscar un documento cuyo nombre ya no recuerda.
  \item \textbf{Somos una empresa 100\% virtual.} No hay pasillo ni café donde
    el conocimiento se transmita de forma informal, así que lo que no está
    escrito prácticamente no existe. La documentación compartida por
    departamento no es un complemento: es el medio principal por el que el
    conocimiento circula.
  \item \textbf{La responsabilidad está repartida.} Como cada departamento
    mantiene lo suyo, la información la actualiza quien la conoce mejor, y no
    un equipo central que se convierte en cuello de botella.
\end{itemize}

La condición para que funcione es la disciplina de actualización. NotebookLM
responde con lo que hay en Drive, de modo que un documento desactualizado
produce una respuesta desactualizada, eso sí, bien citada. Como advierten
\textcite{davenport1998}, un repositorio de conocimiento sólo aporta valor
mientras la organización cultive el hábito de alimentarlo; la herramienta
facilita el acceso, pero no sustituye ese hábito.

\printbibliography[title={Referencias}]

\end{document}
